\documentclass[a4paper,12pt]{article}
\usepackage{amsmath,amssymb}
\usepackage[russian]{babel}
\usepackage[utf8]{inputenc}
\usepackage[T2A]{fontenc}
\usepackage{geometry}
\geometry{margin=1.5cm}

\begin{document}

\section*{Контрольные вопросы и ответы}

\subsection*{1. Чем вызвана необходимость аппроксимирования табличных функций?}

Аппроксимирование табличных функций нужно тогда, когда зависимость известна только в узлах таблицы, а требуется вычислить значение между ними, упростить дальнейшие расчёты или заменить сложную зависимость более простой аналитической формулой. В экспериментальных задачах это позволяет получить удобную математическую модель по результатам измерений.

\subsection*{2. Чем отличается аппроксимация от интерполяции?}

При интерполяции функция строится так, чтобы проходить точно через все узлы таблицы: $\varphi(x_i)=y_i$. При аппроксимации точное совпадение не требуется; наоборот, ищется функция, которая в среднем как можно ближе к исходным данным, то есть $\varphi(x_i)\approx y_i$. Поэтому аппроксимация обычно устойчивее к шуму и ошибкам измерений.

\subsection*{3. Сформулируйте задачу аппроксимации.}

Дана система точек $(x_i,y_i)$, $i=1,\dots,n$. Требуется подобрать функцию $\varphi(x)$ из выбранного класса так, чтобы отклонение между $y_i$ и $\varphi(x_i)$ было минимальным по некоторому критерию. В методе наименьших квадратов таким критерием является сумма квадратов ошибок
\[
S = \sum_{i=1}^{n}\bigl(\varphi(x_i)-y_i\bigr)^2.
\]

\subsection*{4. Как выбирается вид аппроксимирующего уравнения?}

Вид эмпирической зависимости выбирают по характеру экспериментальных точек, по смыслу задачи и по простоте дальнейшего использования модели. На практике сначала рассматривают простые и часто встречающиеся зависимости: линейную, полиномиальную, экспоненциальную, логарифмическую и степенную. После вычисления параметров сравнивают качество приближения и оставляют ту модель, у которой ошибка меньше.

\subsection*{5. Объясните суть метода наименьших квадратов (МНК).}

Суть МНК состоит в том, чтобы выбрать параметры аппроксимирующей функции так, чтобы сумма квадратов отклонений между табличными значениями и значениями модели была минимальной. Это даёт наилучшее приближение в смысле среднего квадрата ошибки:
\[
S = \sum_{i=1}^{n}\bigl(\varphi(x_i)-y_i\bigr)^2 \to \min.
\]
После раскрытия скобок и дифференцирования по неизвестным коэффициентам получается система линейных уравнений, из которой находятся параметры модели.

\subsection*{6. Что такое мера отклонения и как ее вычислить?}

Мера отклонения --- это численная характеристика того, насколько значения аппроксимирующей функции отличаются от исходных табличных данных. В МНК в качестве меры отклонения используется сумма квадратов ошибок:
\[
S = \sum_{i=1}^{n} \varepsilon_i^2 = \sum_{i=1}^{n}\bigl(\varphi(x_i)-y_i\bigr)^2.
\]
Чем меньше $S$, тем точнее модель описывает исходные данные.

\subsection*{7. К решению какой задачи сводится МНК?}

МНК сводится к задаче минимизации функции многих переменных $S(a_0,a_1,\dots,a_m)$ по параметрам модели. Для линейных по коэффициентам функций эта задача эквивалентна решению системы линейных алгебраических уравнений (системы нормальных уравнений).

\subsection*{8. Сформулируйте задачу полиномиальной аппроксимации МНК.}

Требуется для заданных точек $(x_i,y_i)$ найти полином степени $m$:
\[
P_m(x)=a_0+a_1x+\dots+a_mx^m,
\]
такой, чтобы сумма квадратов отклонений
\[
S=\sum_{i=1}^{n}\bigl(P_m(x_i)-y_i\bigr)^2
\]
была минимальной. Неизвестные коэффициенты $a_0,\dots,a_m$ находятся из системы нормальных уравнений.

\subsection*{9. Что такое линейная и квадратичная аппроксимация?}

Линейная аппроксимация --- это приближение данных функцией первой степени:
\[
\varphi(x)=ax+b.
\]
Квадратичная аппроксимация --- приближение полиномом второй степени:
\[
\varphi(x)=a_0+a_1x+a_2x^2.
\]
Квадратичная модель более гибкая, так как может описывать кривизну данных.

\subsection*{10. Приведите графическую интерпретацию линейной и квадратичной аппроксимаций.}

Графически линейная аппроксимация --- это прямая, проходящая «в среднем» максимально близко к облаку экспериментальных точек. Квадратичная аппроксимация --- парабола, которая за счёт дополнительного коэффициента лучше повторяет изгиб данных. В обоих случаях подбирается кривая, минимизирующая суммарные квадраты вертикальных расстояний от точек до кривой.

\subsection*{11. Что такое среднеквадратическое отклонение?}

Среднеквадратическое отклонение аппроксимации (в контексте лабораторной) --- это корень из среднего значения квадратов ошибок:
\[
\sigma = \sqrt{\frac{1}{n}\sum_{i=1}^{n}\bigl(\varphi(x_i)-y_i\bigr)^2} = \sqrt{\frac{S}{n}}.
\]
Оно показывает «типичный» размер ошибки в единицах исходной величины $y$.

\subsection*{12. Как выполняется аппроксимации данных неполиномиальными функциями?}

Для неполиномиальных моделей (экспоненциальной, степенной, логарифмической) часто выполняют линеаризацию заменой переменных, например:
\begin{itemize}
	\item $y=ae^{bx} \Rightarrow \ln y = \ln a + bx$;
	\item $y=ax^b \Rightarrow \ln y = \ln a + b\ln x$;
	\item $y=a\ln x + b$ уже линейна по параметрам $a,b$.
\end{itemize}
После линеаризации коэффициенты определяются тем же МНК. При этом нужно учитывать ограничения области определения (например, $x>0$, $y>0$ для логарифмирования).

\subsection*{13. Как оценить качество полученной аппроксимации?}

Качество оценивают по нескольким критериям:
\begin{itemize}
	\item по сумме квадратов ошибок $S$;
	\item по среднеквадратическому отклонению $\sigma$;
	\item по коэффициенту детерминации $R^2$ (чем ближе к 1, тем лучше);
	\item по графику остатков $\varepsilon_i$ и визуальному совпадению кривой с данными.
\end{itemize}

\subsection*{14. Как выбирается наилучшая аппроксимирующая функция?}

Из нескольких допустимых моделей выбирают ту, для которой ошибка минимальна (обычно по $\sigma$ или $S$) и при этом качество объяснения данных ($R^2$) максимально. Дополнительно учитывают физический смысл модели и отсутствие систематического характера в остатках.

\subsection*{15. Корректно ли применять аппроксимирующие уравнения за пределами исследуемого диапазона?}

В общем случае некорректно или рискованно. Вне диапазона исходных данных выполняется экстраполяция, и ошибка может резко расти, особенно для высоких степеней полинома и неполиномиальных моделей. Поэтому применять модель за пределами исследованного интервала можно только с осторожностью и при наличии дополнительных обоснований.

\end{document}

